% !PRINTED 615
\weberpage{615}

%<<CONT
is formed. Then $c_{1,1}$ is the least positive integral rational number which is divisible by $\varphi_1 \varphi_2$, and this must in any case be a divisor of $a_{1,1} b_{1,1}$. If we have
\[
a_{1,1} b_{1,1} = k\, c_{1,1},
\]
where $k$ is a positive integral rational number, then, since $N_a(\varphi_1)\, N_a(\varphi_2) = N_a(\varphi_3)$ must hold,
\[
k\, a_{2,2} b_{2,2} = c_{2,2}.
\]
The last number $c_{1,2}$ is then to be determined from the congruence
\weq{4}{c_{1,2} + c_{2,2}\,\Theta \equiv 0 \ (\mathrm{mod}\ \varphi_1 \varphi_2).}

If it is only a matter of finding a representative of the composed class, one proceeds most simply as follows.

If $\varphi_1$ is assumed arbitrarily, one can (by \S.\,152, 6.) choose $\varphi_2$ in its class so that it becomes relatively prime to $a_{1,1}$. Then $\varphi_1$ is also relatively prime to $\varphi_2$, and $b_{1,1}$ relatively prime to $a_{1,1}$.

For if $a_{1,1}$ and $b_{1,1}$ have a greatest common divisor $d$, then this is also relatively prime to $\varphi_2$, and hence $b_{1,1} : d$ is divisible by $\varphi_2$; but if $d > 1$, this contradicts the definition of $b_{1,1}$. In the same way it also follows conversely that, if $a_{1,1}$ and $b_{1,1}$ are relatively prime, then $\varphi_2$ is relatively prime to $a_{1,1}$. For if none of the rational prime factors of $b_{1,1}$ divides $a_{1,1}$, then none of the prime factors dividing $\varphi_2$ can divide $a_{1,1}$ either.

In this case we now have:
\weq{5}{a_{1,1} b_{1,1} = c_{1,1}, \quad a_{2,2} b_{2,2} = c_{2,2}.}
The congruence (4) then becomes
\weq{6}{c_{1,2} + a_{2,2} b_{2,2}\,\Theta \equiv 0 \ (\mathrm{mod}\ \varphi_1 \varphi_2),}
and since
\[
a_{1,2} + a_{2,2}\,\Theta \equiv 0 \ (\mathrm{mod}\ \varphi_1), \quad b_{1,2} + b_{2,2}\,\Theta \equiv 0 \ (\mathrm{mod}\ \varphi_2)
\]
holds, (6) splits into the two congruences:
\[
\begin{aligned}
c_{1,2} &\equiv b_{2,2} a_{1,2} \ (\mathrm{mod}\ \varphi_1) \\
&\equiv a_{2,2} b_{1,2} \ (\mathrm{mod}\ \varphi_2).
\end{aligned}
\]
But by \S.\,138, 13. these congruences are equivalent to the ordinary numerical congruences
\weq{7}{
\begin{aligned}
c_{1,2} &\equiv b_{2,2} a_{1,2} \ (\mathrm{mod}\ a_{1,1}) \\
&\equiv a_{2,2} b_{1,2} \ (\mathrm{mod}\ b_{1,1}),
\end{aligned}}
by which $c_{1,2}$ is determined with respect to the modulus $c_{1,1}$.
